\title{Towards a structurally resolved human protein interaction network}
\begin{document}
\maketitle

\begin{abstract}
Nature Structural \& Molecular Biology | Volume 30 | February 2023 | 216–225 216 nature structural \& molecular biology Article https://doi.org/10.1038/s41594-022-00910-8 Towards a structurally resolved human protein interaction network David F. Burke 1,9, Patrick Bryant 2,3,9, Inigo Barrio-Hernandez 1,9, Danish Memon 1,9, Gabriele Pozzati 2,3,9, Aditi Shenoy 2,3, Wensi Zhu2,3, Alistair S. Dunham 1, Pascal Albanese4,5, Andrew Keller6, Richard A. Scheltema 4,5, James E. Bruce 6, Alexander Leitner 7, Petras Kundrotas 2,3,8 , Pedro Beltrao 1,7 \& Arne Elofsson 2,3 Cellular functions are governed by molecular machines that assemble through protein-protein interactions. Their atomic details are critical to studying their molecular mechanisms. However, fewer than 5\% of hundreds of thousands of human protein interactions have been structurally characterized. Here we test the potential and limitations of recent progress in deep-learning methods using AlphaFold2 to predict structures for 65,484 human protein interactions. We show that experiments can orthogonally confirm higher-confidence models. We identify 3,137 high-confidence models, of which 1,371 have no homology to a known structure. We identify interface residues harboring disease mutations, suggesting potential mechanisms for pathogenic variants. Groups of interface phosphorylation sites show patterns of co-regulation across conditions, suggestive of coordinated tuning of multiple protein interactions as signaling responses. Fin
\end{abstract}

\section{Recovered Text}
Nature Structural \& Molecular Biology | Volume 30 | February 2023 | 216–225
216
nature structural \& molecular biology
Article
https://doi.org/10.1038/s41594-022-00910-8
Towards a structurally resolved human
protein interaction network
David F. Burke
  1,9, Patrick Bryant
  2,3,9, Inigo Barrio-Hernandez
  1,9,
Danish Memon
  1,9, Gabriele Pozzati
  2,3,9, Aditi Shenoy
  2,3, Wensi Zhu2,3,
Alistair S. Dunham
  1, Pascal Albanese4,5, Andrew Keller6,
Richard A. Scheltema
  4,5, James E. Bruce
  6, Alexander Leitner
  7,
Petras Kundrotas
  2,3,8
, Pedro Beltrao
  1,7
 \& Arne Elofsson
  2,3
Cellular functions are governed by molecular machines that assemble
through protein-protein interactions. Their atomic details are critical to
studying their molecular mechanisms. However, fewer than 5\% of hundreds
of thousands of human protein interactions have been structurally
characterized. Here we test the potential and limitations of recent progress
in deep-learning methods using AlphaFold2 to predict structures for 65,484
human protein interactions. We show that experiments can orthogonally
confirm higher-confidence models. We identify 3,137 high-confidence
models, of which 1,371 have no homology to a known structure. We identify
interface residues harboring disease mutations, suggesting potential
mechanisms for pathogenic variants. Groups of interface phosphorylation
sites show patterns of co-regulation across conditions, suggestive of
coordinated tuning of multiple protein interactions as signaling responses.
Finally, we provide examples of how the predicted binary complexes can
be used to build larger assemblies helping to expand our understanding of
human cell biology.
Proteins are key cellular effectors determining most cellular processes.
These rarely act in isolation, but instead, the coordination of the diver-
sity of processes arises from the interaction among multiple proteins
and other biomolecules. The characterization of protein-protein inter-
actions (PPIs) is crucial for understanding which groups of proteins
form functional units and underlies the study of the biology of the cell.
Diverse experimental and computational approaches have been devel-
oped to determine the PPI network of the cell (that is, the interactome),
with hundreds of thousands of human protein interactions determined
to date1–3. Protein interactions vary from transient interactions that
regulate an enzyme to permanent interactions in molecular machines.
The structural characterization of the human interactome has
lagged behind, with experimental and homology models currently
covering an estimated 15 protein interactions4,5. The structural char-
acterization of protein complexes is a critical step in understand-
ing the mechanisms of protein function, and in studying the impact
of mutations4,6–8 and the regulation of cellular processes via the
post-translational tuning of binding affinities9–12.
Received: 11 February 2022
Accepted: 14 December 2022
Published online: 23 January 2023
 Check for updates
1European Molecular Biology Laboratory, European Bioinformatics Institute (EMBL-EBI), Cambridge, UK. 2Science for Life Laboratory, Stockholm
University, Solna, Sweden. 3Department of Biochemistry and Biophysics, Stockholm University, Stockholm, Sweden. 4Biomolecular Mass Spectrometry
and Proteomics, Bijvoet Center for Biomolecular Research and Utrecht Institute of Pharmaceutical Sciences, Utrecht University, Utrecht, The Netherlands.
5Netherlands Proteomics Center, Utrecht, The Netherlands. 6Department of Genome Sciences, University of Washington Seattle, Seattle, WA, USA.
7Department of Biology, Institute of Molecular Systems Biology, ETH Zurich, Zurich, Switzerland. 8Center for Computational Biology, The University of
Kansas, Lawrence, KS, USA. 9These authors contributed equally: David F. Burke, Patrick Bryant, Inigo Barrio-Hernandez, Danish Memon and Gabriele
Pozzati.
 e-mail: pkundro@ku.edu; pbeltrao@ebi.ac.uk; arne@bioinfo.se
Nature Structural \& Molecular Biology | Volume 30 | February 2023 | 216–225
217
Article
https://doi.org/10.1038/s41594-022-00910-8
the random set of models would be considered confident predictions
at this cut-off. In Fig. 1c we show examples of predicted structures
aligned to experimental or homology models, showing how the pre-
dictions and the confidence score relate to the observed alignments.
For the majority of these cases, even with lower-confidence values, the
interaction interface is generally in good agreement, except for the
interaction between subunits of the proteasome 26S complex, ATP-
pase domain 2 (PSMC2) and non-ATPase domain 11 (PSMD11). It can be
noted that several of the models in Fig. 1c are parts of large complexes:
PRDX2–PRDX3: members of the peroxiredoxin family of antioxidant
enzymes; RFC2–RFC5: subunits of heteropentameric Replication fac-
tor C (RF-C); YWHAB–YWHAG: parts of the 14-3-3 family of proteins
tyrosine 3-monooxygenase/tryptophan 5-monooxygenase activation
proteins beta (YWHAB) and gamma (YWHAG); and RPL9–RPL18A:
ribosomal proteins L9 (RPL9) and L18a (RPL18A). This shows that Alpha-
Fold2 can predict the structures of directly interacting protein pairs
present in large complexes.
Features impacting prediction confidence
As shown in Fig. 1a, protein pairs present in the Protein Data Bank (PDB)
are enriched in high-scoring models compared with pairs in HuRI and
Hu.MAP. There could exist several possible explanations for this, such
as the inability of AlphaFold2 to identify transient or indirect interac-
tions. Nevertheless, it is also possible that the two high-throughput
datasets contain noninteracting pairs. Therefore, to understand this
difference better, we first studied an additional dataset created from
large (>10 chains) heteromeric protein complexes.
The set of large complexes consists of 12 large heteromeric protein
complexes, and all (nonidentical) pairs of protein chains in each com-
plex were docked with each other. These pairs can be divided into the
ones with direct interaction and those that do not interact directly. Here,
we used a definition of more than 20 contacts of less than 8 Å between
Calphas to exclude small interaction interfaces. When a complex con-
tained multiple copies of identical chains, all interactions were included
to allow for alternative interactions between the chains. The difference
in pDockQ scores between the direct and indirect interacting pairs is
striking, where only 6\% of the indirect pairs have a pDockQ score > 0.5
compared with 38\% of the directly interacting pairs (Fig. 2a). This shows
that directly interacting pairs often can be predicted even when they
are part of large complexes, in contrast to indirectly interacting pairs.
hu.MAP has many more high-confidence predictions than HuRI,
which is based on Y2H experiments. To further understand this differ-
ence, we first analyzed a subset of all protein pairs from the CORUM23
database, the best manually curated database of mammalian protein
complexes, and predicted the interaction of all pairs in the same com-
plex. The average pDockQ score of CORUM is slightly higher than for
hu.MAP, but the number of high-quality predictions is similar (16\% ver-
sus 19\%), indicating that the different databases of protein complexes
have a similar fraction of high-confidence predictions and that HuRI
is the outlier (Fig. 2b).
It is unlikely that the Y2H in HuRI data should contain a large set
of indirect interactions, as only two human proteins are expressed
in the same cell. Therefore, there must be another reason for the few
high-confidence predictions. We examined the properties of the pairs
present in the two datasets. Here, it can be seen that HuRI proteins
differ from the hu.MAP (and other datasets) in two ways. HuRI protein
pairs contain more intrinsic disorder (Fig. 2c) and have fewer efficient
sequences (meff) in their multiple sequence alignments (MSAs) (Fig. 2d).
In these figures it can also be seen that the pDockQ values tend to
increase with less disorder and more sequences in the alignments,
although it is clearly not an absolute relationship. Further, protein pairs
in HuRI are less likely to be found in the same subcellular compartment
(Fig. 2e), and have similar coexpression profiles (Fig. 2f). Considering
all this, it is likely that many protein interactions in HuRI are transient
and that AlphaFold2 cannot reliably predict such interactions.
Computational approaches for predicting the structures of inter-
acting protein pairs are primarily based on identifying structural simi-
larity for pairs of proteins against experimentally determined protein
complexes4,6,13,14. The Interactome3D (refs. 4,14) repository currently
lists 7,625 predicted models based on homology of domains, a number
similar to the 8,359 pairs listed having an experimentally determined
model. In addition, co-evolution-based information has been used
to predict protein interactions and to guide structural docking for
bacterial proteins15. Recently, neural network-based approaches have
demonstrated the ability to accurately predict the structures of indi-
vidual proteins16,17 and protein complexes16,18–21. These approaches
can correctly predict the structures of up to 60\% of dimers18, and have
been used to predict structures of 1,506 Saccharomyces cerevisiae pro-
tein interactions22. However, the application of these neural network
models for the large-scale prediction of human complex structures
has not been tested yet.
Here, we assess the possibilities and limitations of applying Alpha-
Fold2 to modeling human protein interactions on a large scale. We
predicted the complex structures for two sets of human interactions
obtained using different experimental methods, comprising 65,484
unique human interactions. We show that it is possible to rank the mod-
els according to confidence, with 3,137 predicted structures ranked as
highly confident. Further, we show that the higher-confidence predic-
tions are enriched among those supported by a combination of experi-
mental methods. We showcase the value of a structurally resolved
interactome by studying disease mutations and phosphorylation of
interface residues. Finally, we provide some indication that binary
complexes can be used to build higher-order assemblies.
Structure prediction of human protein
interactions
We selected experimentally identified human protein interactions
from the Human Reference Interactome (HuRI)2 and the Human Protein
Complex Map (hu.MAP v.2.0)3. HuRI comprises protein interactions
determined by yeast two-hybrid (Y2H) screening2 from which we mod-
eled 55,586 pairs. From hu.MAP we selected 10,207 high-quality PPIs3.
While HuRI is more likely to be enriched for direct protein interactions,
including transient partners, the hu.MAP set is more likely to reflect sta-
ble protein interactions, including members of the same complex that
may not be interacting directly. The overlap between the two datasets
is small (309 pairs), and a comparison with two large-scale compendi-
ums of structural models4 indicates that 62,019 of the combined pairs
do not have experimental models nor can they be modeled easily by
homology, suggesting a large potential gain in structural knowledge.
We predicted the structure of 65,484 nonredundant pairs using the
FoldDock pipeline18, based on AlphaFold2 (ref. 17). As in the FoldDock
pipeline, we combined size and the predicted local Distance Differ-
ence Test (plDDT) scores of the interface into a single score to predict
the DockQ score of a complex, dubbed pDockQ (Methods), which can
rank models by confidence. We tested pDockQ score by comparing the
predicted models with 1,465 experimental models, of which 742 (50\%)
were correct (DockQ > 0.23). For predictions with pDockQ > 0.23, 70\%
(671 of 955) are well modeled, and for pDockQ > 0.5, 80\% (521 of 651).
We show in Fig. 1a the distribution of pDockQ for the predicted and
random protein interactions, and provide data for all models in Sup-
plementary Table 1. The pDockQ of known interacting proteins tends
to be higher than for the random set, with the predictions for hu.MAP
showing on average higher confidence than the HuRI set. Additionally,
when selecting hu.MAP interactions also supported by Y2H or crosslink
data (crosslinking) results in even higher-confidence values (Fig. 1a).
This suggests that high-confidence models are enriched for protein
interactions supported by the two types of methods associated with
high affinity and direct interactions. We identified 3,137 structures
(Fig. 1b) as high-confidence models (pDockQ > 0.5). The number of
structures increased to 10,061 if a cut-off of 0.23 was used. Only 0.3\% of
Nature Structural \& Molecular Biology | Volume 30 | February 2023 | 216–225
218
Article
https://doi.org/10.1038/s41594-022-00910-8
Crosslinking support for predicted complex
structures
Chemical crosslinking followed by mass spectrometry is an approach
which can be used to identify reactive residues (usually lysines) that
are in proximity, as constrained by the geometry of the crosslink agent
used. The identification of such residues across a pair of proteins can
help define the likely protein interface. To determine if the predicted
complex structures agree with such orthogonal spatial constraints, we
obtained a compilation of crosslinks for pairs of residues across 528
protein pairs with predicted models (Fig. 3a, Supplementary Table 1
and Methods). In total, 51\% of the models had one or more crosslinks
at a distance below the expected maximal distance possible (Fig. 3a).
Restricting the predicted models to higher confidence by the pDockQ
score increased the fraction of complexes with acceptable crosslinks,
reaching 75\% for pDockQ scores greater than 0.5 (Fig. 3a). This result
is in line with the benchmark results above.
In total, we have identified 479 crosslinks providing supporting
evidence for 171 predicted complex structures with pDockQ > 0.5.
Of these, 41 correspond to complex structures with no experimen-
tal structure or homology models, from which we selected some to
illustrate (Fig. 3b–e). Figure 3b shows the AlphaFold2 (AF2) model
for the full length of the ERLIN1/ERLIN2 complex, which mediates
the endoplasmic reticulum-associated degradation (ERAD) of ino-
sitol 1,4,5-trisphosphate receptors (IP3Rs). AlphaFold2 predicts a
globular domain (1–190) followed by an extended helical region with
a kink around amino acid position 280. Unlike the model in Interac-
tome3D, the paralogous proteins are stacked side-by-side with the
hydrophobic face of the helices buried and the hydrophilic face
(mainly Lys) exposed to solvent. A crosslink between the C-terminal
residues K275 (ERLIN1) and K287 is predicted to bridge a distance of
18 Å, supporting the predicted model. In Fig. 3c we show the model
for proteins IMMT and CHCHD3, components of the mitochondrial
inner membrane MICOS complex. AlphaFold2 predicts a globular
helical domain at the C-terminal end of IMMT (550–750) to interact
with the C-terminal end of CHCHD3 (150–225). This is supported by
data of three crosslinks: between K173 (CHCD3) and K565 (IMMT), and
K203 (CHCD3) to both K714 and K726 of IMMT. Figure 3d shows the
complex of transfer RNA-guanine-N(7)-methyltransferase (METTL)
with its noncatalytic subunit (WDR4). The structure of WDR4 has not
yet been solved experimentally but contains WD40 repeats, which
are expected to form a β-propeller domain, as predicted here. The
METTL domain is predicted to interact with the side of the WDR40,
away from the ligand-binding pore. This orientation is supported by
a crosslink between K122 (WDR4) and K143 (METTL) (18 Å). Finally, in
Fig. 3e we show the predicted complex structure for the heterogene-
ous nuclear ribonucleoprotein C (HNRNPC) and the RNA-binding
protein, RALY. Two regions in both proteins are predicted with high
confidence (plDDT > 70), with the lower-confidence regions not shown.
The N-terminal domain in HNRNPC (16–85) is predicted to interact
with the N-terminal domain of RALY (1–100). A long helix in HNRNPC
(185–233) is predicted to interact with a helix in RALY (169–228). This
interhelix interface is supported by crosslinking data for three pairs of
lysines at either end of the helices (189 → 222; 229 → 179; and 232 → 183).
Disease-associated missense mutations at
interfaces
Missense mutations associated with human diseases can alter pro-
tein function via diverse mechanisms, including disrupting protein
stability, allosterically modulating enzyme activity and altering PPIs.
Structural models can allow the rationalization of possible mechanisms
DockQ = 0.74
pDockQ = 0.68
DockQ = 0.83
hu.MAP 2.0
HuRI
Structure
Y2H
Crosslinking
Random
pDockQ
Density
16
14
12
10
8
6
4
2
0
–0.25
0
0.25
0.50
0.75
1.00
pDockQ = 0.73
DockQ = 0.69
pDockQ = 0.69
DockQ = 0.49
pDockQ = 0.08
DockQ = 0.01
pDockQ = 0.03
1,404
8,494
1,536
53,741
HuRI
hu.MAP
Predicted
0
197
112
YWHAB
RFC2
RPL9
PSMC2
PSMD11
YWHAG
RFC5
RPL18A
PRDX2
PRDX3
a
b
c
Fig. 1 | Application of AlphaFold2 complex predictions to a large dataset of
human PPIs. a, Distribution of model confidence score (pDockQ) for predicted
structures from two large human protein interaction datasets (hu.MAP and
HuRI), compared with confidence metrics from 2,000 random pairs of proteins.
The hu.MAP dataset was further subsetted to those that have support from Y2H
(‘Y2H’) or crosslink data (‘Crosslinking’), or correspond to pairs with available
experimental or homology modeling information (‘Structure’). b, Number of
protein interactions with models built from both datasets and those that we
consider being of high confidence (‘Predicted’), corresponding to those with
pDockQ > 0.5. c, Examples of predicted models (orange and green) overlapped
with the corresponding experimental models (gray) and the observed (DockQ)
or predicted (pDockQ) quality of the models.
Nature Structural \& Molecular Biology | Volume 30 | February 2023 | 216–225
219
Article
https://doi.org/10.1038/s41594-022-00910-8
of interface disease mutations. To determine the usefulness of the
predicted structures, we compiled a set of mutations located at
interface residues that were previously experimentally tested for
the impact on the corresponding interaction24. We then performed
in silico predictions of changes in binding affinity upon mutations
using FoldX25 and observed that mutations known to disrupt the
interactions are predicted to have a strong destabilization of binding
compared with mutations known not to have an effect (Fig. 4a and
Supplementary Table 2). Very high confidence (plDDT > 90) of the
mutated residues led to more substantial discrimination between
mutations known and not known to disrupt the complex formation
(Fig. 4a), indicating that only very accurate models are useful when
using the FoldX forcefield for estimating the impact of binding affinity
of mutations.
Next, we mapped human disease (from ClinVar) and cancer
mutations (from The Cancer Genome Atlas) to the interface residues
defined by the set of high-confidence protein complex predictions
(pDockQ > 0.5) (Supplementary Table 1). The hu.MAP and HuRI con-
fident predictions identified 280 interfaces carrying pathogenic
mutations and 602 interfaces corresponding to the top 25\% of recur-
rently mutated interfaces in cancer, defined as the highest number
of mutations per interface position (Fig. 4b and Methods). We find a
strong enrichment in pathogenic versus benign mutations at interface
residues relative to the rest of the protein (2.3-fold enrichment, P value
2.7 × 10−31).
We illustrate in Fig. 4c examples of protein network clusters with
interface disease mutations across a range of biological functions.
For example, interface mutations in chromatin remodeling, including
members of SWI/SNF complex (SMARCD1, SMARCD2, SMARCD3), and
several transcription factors related to development (for example,
TCF3, TCF4, LMO1 and LMO2).
We selected examples of interfaces with disease mutations and no
previous experimental data or homology to available models (Fig. 4d–g).
Figure 4d shows the interface of WDR4-METTL1, which has support-
ing crosslink information described above. WDR4 has two annotated
pathogenic variants at this interface, linked with Galloway-Mowat Syn-
drome 6, with the highlighted R170 participating in interactions with
a negatively charged residue of METTL1. Figure 4e shows an example
of an interface with 32 recorded interface mutations in cancer for both
proteins, including the highlighted arginines in LDOC1, which form
electrostatic interactions with the opposite chain. TWIST1 has several
annotated pathogenic mutations, including L149R and L159H, which
a
HuRI: 0.06
8
7
6
5
Density
pDockQ
pDockQ
pDockQ
FracDiso
4
3
2
1
0
–0.25
0
0.25
0.50
0.75
1.00
log(meff)
SubcellularDegree
0
1
2
3
Different
Organell
SubOrg
SubSub
4
0
0.2
0.4
0.6
0.8
0
0.2
0.4
0.6
0.8
1.0
gtex
0
0.2
0.4
0.6
0.8
1.0
8
7
6
5
Density
4
3
2
1
0
0
0.2
0.4
0.6
0.8
pDockQ
0
0.2
0.4
0.6
0.8
pDockQ
0
0.2
0.4
0.6
0.8
pDockQ
0
0.4
0.2
–0.2
0.6
0.8
1.0
hu.MAP: 0.16
Direct: 0.38
Indirect: 0.06
HuRI: 0.055
hu.MAP: 0.159
CORUM: 0.19
HuRI
hu.MAP
Dataset
HuRI
hu.MAP
Dataset
HuRI
hu.MAP
Dataset
HuRI
hu.MAP
Dataset
b
c
d
e
f
Fig. 2 | Protein and interaction features impacting on prediction confidence:
analysis of different datasets. In all subfigures, proteins from HuRI in green,
hu.MAP gray, CORUM orange and from large PDB complexes blue. a, pDockQ
values of directly and indirectly interacting proteins from the same complex
(blue); for comparison, HuRI and hu.MAP data are shown with thin lines.
b, pDockQ values of CORUM (orange), HuRI (green) and hu.MAP (gray) datasets.
c, Fraction of residues predicted to be disordered (pLDDT < 0.5) shows that
protein pairs in HuRI are enriched in disorder. d, Proteins in HuRI have fewer
sequences in the paired MSAs as measured by the mean number of efficient
sequences in the MSA (meff). e, Proteins that share subcellular localization
(solid lines) are enriched in high pDockQ scores in all three datasets. f, Only
protein pairs in hu.MAP are coexpressed according to STRING, using similarity in
Genotype-Tissue Expression (gtex), and coexpressed pairs are enriched in pairs
with high pDockQ scores.
Nature Structural \& Molecular Biology | Volume 30 | February 2023 | 216–225
220
Article
https://doi.org/10.1038/s41594-022-00910-8
are at residues buried in the interface (Fig. 4f). In particular, the L149R
mutation, associated with Saethre–Chotzen syndrome, would strongly
disrupt packing. The R118G mutation would disrupt the interaction with
residue F22 mainchain O in TCF4. In RAD51D we found the mutation
R266C (Breast-ovarian cancer, familial), which interacts across the
interface with XRCC2 (Fig. 4g) and paralogous genes involved in the
repair of DNA double-strand breaks by homologous recombination.
Interestingly, we also found mutations at R239, to Trp/Gln/Gly, associ-
ated with Breast-ovarian cancer which interacts with Tyr119 in XRCC2,
which itself is also annotated as having mutations linked to hereditary
cancer-predisposing syndrome.
Phospho-regulation of protein complex
interfaces
Protein phosphorylation can regulate protein interactions by modulat-
ing the binding affinity via the change in size and charge of the modified
residue. Over 100,000 experimental human phosphorylation sites
have been determined to date26,27, but only 5–10\% of these have a known
function28. Mapping phosphorylation site positions to protein inter-
faces can generate mechanistic hypotheses for their functional roles in
controlling protein interactions. We used a recent characterization of
the human phosphoproteome26 to identify 4,145 unique phosphosites
at interface residues among the highly confident models. The aver-
age functional importance, defined by the functional score described
earlier26, was generally higher than random for phosphorylation sites
at interfaces (Fig. 5a), and we found some enrichment for targets of
multiple kinases, including tyrosine kinases (ERBB2, AXL, ABL2, FER)
(Fig. 5b). This suggests that some interfaces may be under coordinated
regulation by specific kinases and conditions.
To identify potentially co-regulated interfaces, we collected meas-
urements of changes in phosphorylation levels across a large panel of
over 200 conditions29. We retained 260 phosphosites that had a sig-
nificant regulation in three conditions and then computed all-by-all
pairwise correlations in phosphosite fold changes across conditions
(Supplementary Table 1). We clustered these phosphosites by their
profile of correlations (Fig. 5c), identifying 16 groups of co-regulated
interface phosphorylation sites (Fig. 5c and Supplementary Table 3). For
each group of phosphosites, we identified the conditions where these
have the strongest up- or down-regulation (Supplementary Fig. 1) and
plotted a subset of conditions in Fig. 5d. We also performed a gene ontol-
ogy enrichment analysis for each group of co-regulated phosphosites,
including both proteins of the modified interfaces, to search for common
biological functions (Fig. 5e and Supplementary Table 4). Here, one-sided
hyper-geometric tests were used for statistical analysis. For example,
we observed a cluster of interface phosphosites in proteins related to
intermediate filaments (cluster 7) which show strong regulation pat-
terns along the cell cycle, downregulated in S-phase and up-regulated
in G1 and mitosis. Phosphosites in cluster 1 (cell cycle G1-S phase transi-
tion) show the opposite trends, with up-regulation in late S-phase and
down-regulation in G1 and mitosis. Some clusters show regulation under
specific kinase inhibition, which may provide novel hypotheses for kinase
regulation of specific processes. For example, phosphosites in cluster
9 (regulation of chromosome assembly) tend to be up-regulated after
inhibition of ROCK and up-regulated after inhibition of mTOR.
While not all phosphosites at interfaces are likely to regulate the
binding affinity, this analysis provides hypotheses for the potentially
coordinated regulation of multiple proteins by tuning of their interac-
tions after specific perturbations.
0
100 200 300 400 500 600
Xlink count
ERLIN1
ERLIN2
pDockQ = 0.73
Int plDDT = 0.87
Int residues = 400
WDR4
HNRNPC
METTL1
RALY
pDockQ = 0.70
Int plDDT = 0.91
Int residues = 198
pDockQ = 0.60
Int plDDT = 0.86
Int residues = 157
CHCHD3
IMMT
pDockQ = 0.62
Int plDDT = 84
Int residues = 185
a
b
c
d
e
pDockQ<0.23
All pairs
0.23<
pDockQ<0.5
pDockQ>0.5
0
0.2 0.4 0.6 0.8
Ratio (valid Xlink)
All Xlink
Xlink<32Å
Fig. 3 | Crosslink support for predicted complex models. a, The numbers and
ratios of predicted structures having crosslink information for pairs of residues
that bridge the two proteins in the predicted structure, broken down by the
crosslinks that satisfy their expected maximal distance and by the predicted
quality of the model (pDockQ). b–e, Examples of predicted structures of high
confidence, with no previous structural information and supported by at least
one crosslink (indicated with blue line): ERLIN1/ERLIN2 complex (b), IMMT and
CHCHD3, components of the mitochondrial inner membrane MICOS complex
(c), the complex of transfer RNA-guanine-N(7)-methyltransferase (METTL)
with its noncatalytic subunit (WDR4) (d) and the heterogeneous nuclear
ribonucleoprotein C (HNRNPC) and the RNA-binding protein, RALY (e).
Nature Structural \& Molecular Biology | Volume 30 | February 2023 | 216–225
221
Article
https://doi.org/10.1038/s41594-022-00910-8
Higher-order assemblies from binary protein
interactions
Proteins interact with multiple partners either simultaneously, as part
of larger protein complexes, or separated in time and space. This is also
reflected in our structurally characterized network, where proteins can
be found in groups, as illustrated in a global network view of the pro-
tein interactions with confident models (Fig. 6, Supplementary Fig. 2
and Supplementary Data 1). One key benefit of structurally character-
izing an interaction network is the identification of shared interfaces
for multiple interactors. As an example, we highlight GDI1 (RabGDP
dissociation inhibitor alpha) which interacts with multiple Rab pro-
teins, regulating their activity by inhibiting the dissociation of GDP.
The predicted complex structures for these interactions show how
these share the same interface and therefore cannot co-occur. Other
clusters in the network suggest that the proteins form larger protein
complex assemblies with many-to-many interactions. As the use of
AlphaFold2 for predicting larger complex assemblies can be limited
by computational requirements, we tested whether the structures for
pairs of proteins could be iteratively structurally aligned. We tested
this procedure on a small set of complexes covered in this network,
with known structures and the number of subunits ranging from five
(RFC complex, TFIIH core complex) to 14 (20S proteasome). We then
aligned an experimentally determined structure with the predicted
models (Fig. 6; gray, experimental model). These examples showcase
the potential and also limitations of this procedure.
The TFIIH core complex is composed of five subunits with 1-to-1
stoichiometry. All subunits can be modeled, with the final complex
generally agreeing (Fig. 6) with a cryoEM structure for these subunits
(PDB:6NMI). The most significant difference to the cryoEM model is
the relative positioning of the ERCC3 subunit. The exact final model
Pathogenic genetic variation
(ClinVar)
Edge shape
YES
NO
Edge color
YES
NO
Gene mutations in cancer
(top 25\%)
5.1 × 10−05
0.053
0.708
0.016
3.2 × 10−05
Very low (pLDDT < 50)
Low (70 > pLDDT > 50)
Confident (90 > pLDDT > 70)
Very high (pLDDT > 90)
Experimental
−2.5
0
2.5 5.0 7.5
∆∆G
Impact
Neutral
Unique PPI count (high confidence)
All pairs
Mutations
present in cancer
0
1,000
2,000
3,000
4,000
3,137
3,022
602
280
Mutations present
in cancer (high ratio)
Human variation associated
with pathogenic phenotypes
PEG10
LDOC1
TCF4
RAD51D
TWIST1
XRCC2
R31
R118
L149
R266
R239
R63
Y119
d
e
f
g
WDR4
METTL1
R164
TCF4
FERD3L
ID2
ID1
MSGN1
MYF5
TAL2
NHLH2
NEUROG2
HES6
TWIST2
TWIST1
NEUROG3
ASCL3
NEUROD1
NEUROD4
HAND1
TCF12
TCF23
TCF3
ASCL4
ID3
Transcription, DNA-templated (development)
SMARCC2
SMARCB1
SMARCD2
SMARCD3
SMARCD1
SMARCC1
HDAC2
Chromatin remodeling
PPP1R3C
SH2D4A
NRAC
PPP1R27
PPP1R11
PPP1CB
PPP1R2C
Dephosphorylation
STX11
SNAP23
STXBP2
STXBP1
STX19
SNAP25
Regulated exocytosis
C11orf49
PRKACB
PRKACG
AVPI1
PRKAR1B
PRKAR2B
PRKAR1A
PRKACA
PRKX
Regulation of kinase activity
LDB2
LMO2
LHX9
LMO4
ISL1
LMO3
LHX8
LMO1
LDB1
LHX3
Transcription, DNA-templated
a
b
c
HAND2
Fig. 4 | Disease mutations at protein complex interface residues. a, Boxplot
showing the distribution of changes in predicted binding affinity (ΔΔG) for
mutations known to have an impact (orange) versus the ones with neutral effect
(green). The boxes represent the first and third quartiles. The upper whisker
extends from the third quartile to the largest value no further than 1.5 × inter-
quartile range (IQR). The lower whisker extends from the first quartile to the
smallest value at most 1.5 × IQR. The statistical significance of the differences
was tested using two-sided Wilcoxon rank sum test with continuity correction.
The total number of data samples was 1,320 (900 neutral, 420 impact)
(Supplementary Table 2). The min, mean and max values for the subsets are
as follows: Experimental Neutral (−0.5, 1.0, 5.1); Experimental Impact (0.1, 3.4,
8.6); AF plddt > 90 Neutral (−0.9, 0.3, 4.7); AF plddt > 90 Impact (−2.7, 1.0, 12.5,);
AF 70 < plddt > 90 Neutral (−2.0, 0.3, 9.0); AF 70 < plddt > 90 Impact (−4.0, 0.5,
7.0); AF 50 < plddt > 70 Neutral (−1.2, 0.1, 2.7); AF 50 < plddt > 70 Impact (−0.2,
0.3, 1.7); AF plddt < 50 Neutral (−2.6, 0.1, 2.0); AF plddt < 50 Impact (−2.6, −0.3,
0.3). b, Unique PPI pairs for high-confidence models (pDockQ > 0.5) in total,
with mutations in cancer, mapped to the interface (all and top 25\% ratios) and
with pathogenic or likely pathogenic clinical variants mapped to the interface.
c, Modules related to relevant biological processes. The color of the edge
represents the presence of cancer mutations in the interface (top 25\% ratio, color
red), and the shape the presence of pathogenic clinical variants (double line).
d–g, Selected relevant structures with no previous structural knowledge showing
clinical variants or mutations in cancer mapped to the interface (mutated
residues in red). The interface of WDR4-METTL1 (d), the interface of LDOC1 (e),
the interface of TWIST1 (f) and the interface of RAD51D (g).
Nature Structural \& Molecular Biology | Volume 30 | February 2023 | 216–225
222
Article
https://doi.org/10.1038/s41594-022-00910-8
obtained can vary depending on the aligned pairs, with multiple possi-
ble final conformations (Supplementary Fig. 3). Figure 6 illustrates the
conformation that best matches the cryoEM model in PDB:6NMI. For
example, for the TFIIH core complex, there is a predicted model where
the complex adopts a more open conformation (as seen in PDB:5OQJ)
and alternative predicted placements of the GTF2H1 subunit.
The RFC complex is also composed of five subunits with 1-to-1
stoichiometry. One iterative alignment of pairwise protein interac-
tions builds a model that includes all five subunits organized similarly
to that observed in the PDB:6VVO cryoEM structure (Fig. 6). In this
predicted model, the subunits RFC2/5/4/3 match the experimentally
observed model well, but there are apparent deviations introduced
by compounding errors in alignment by this iterative process. Indi-
vidual subunits in the cryoEM structure can be aligned to each of the
model subunits well, but then the alignment of the rest of the model
is progressively worse the further away the subunits are positioned
from the aligned subunit. The RFC1 subunit is individually not well
predicted. Further, the RFC3-RFC5 interaction pair is predicted with
high confidence, while, in fact, these do not share a direct contact in
the experimental structure. AlphaFold2 places RFC3 at the RFC5-RFC4
interface, likely due to the structural similarity between RFC3
and RFC4.
a
b
c
d
e
Kinase
Correlation
−1.0
−0.5
0
0.5
1.0
Clusters
1
10
11
12
13
14
15
16
2
3
4
5
6
7
8
9
*
*
*
* *
*
*
*
*
* *
*
*
* *
* *
*
*
*
*
*
*
* *
*
*
*
*
*
* *
* *
*
*
*
*
* *
*
Trichostatin A\_vs\_DMSO
PGE2\_10.00m
Doxorubicin
Bortezomib
MEKi (U0126)
ROCKi (H−1152)
KRAS (G13D)
PI3Ki (GDC−0941)
Early S (Thymidine)
S\_G2 (Thymidine)
Vemurafenib (PLX4032)
Late S (Thymidine)
PP2Ai (Okadaic acid)
G1 (Nocodazole)
AKTi (MK−2206)
mTORi (Torin−1)
Mitosis (Nocodazole)
Clusters
hu.MAP
HuRI
Random
0
0.25
0.50
0.75
1.00
Functional score
−log10(P value)
−10
−5
0
5
10
CL7
CL8
CL6
CL3
CL16
CL4
CL9
CL5
CL11
CL1
CL2
CL12
CL14
1
2
3
5
7
8
9
12
13
14
0
10
20
30
Dna replication initiation
Cell cycle g1 s phase transition
Enzyme binding
Intermediate filament
Intermediate filament cytoskeleton
Structural molecule activity
Regulation of telomere maintenance
Positive regulation of telomere maintenance
Protein stabilization
Atp dependent dna helicase activity
Dna conformation change
Purine ntp dependent helicase activity
Intermediate filament based process
Intermediate filament cytoskeleton
Intermediate filament
Keratin filament
Intermediate filament cytoskeleton
Intermediate filament
Chromosome organization
Swi snf superfamily type complex
Chromatin assembly or disassembly
Threonine type peptidase activity
APC dependent catabolic process
Endopeptidase activity
Protein sumoylation
Nuclear pore
Organelle fission
Keratin filament
Intermediate filament
Intermediate filament cytoskeleton
−log10(P value)
MKNK2
AXL
ABL2
ERBB2
FER
MOS
PKMYT1
CDK7
PRKDC
WEE1
0
1
2
−log10(P value)
Fig. 5 | Co-regulation of phosphorylation sites at interface residues.
a, Distribution of phosphosite functional scores for phosphosites at interface
residues and random phosphosites. The min, mean and max values were as
follows: Random = 0.02, 0.26, 0.98; HuRI = 0.06, 0.37, 0.99; hu.MAP = 0.06,
0.33, 0.99. The boxes represent the first and third quartiles. The upper whisker
extends from the third quartile to the largest value no further than 1.5 × IQR.
The lower whisker extends from the first quartile to the smallest value at most
1.5 × IQR. b, Enrichment of kinase substrates among phosphosites at interface
residues. The P value was derived from an over-representation analysis using
a one-sided hyper-geometric test (N = 7,150). c, Hierarchical clustering of the
pairwise correlation values for changes in phosphosite levels across conditions.
Groups of phosphosites showing high correlation values were defined as
clusters (1 to 16), as indicated in colors along the outside of the clustergram.
d, Degree of regulation of phosphosites from each cluster in a select panel of
conditions, defined by a one-sided Z-test comparing the fold change of the
phosphosites in a cluster compared with the entire distribution of fold changes
in that condition. The result is summarized as the −log(P value), and signed as
positive if the median value is above the background or negative otherwise.
e, Gene ontology enrichment analysis for the proteins with phosphosites
annotated to select clusters.
Nature Structural \& Molecular Biology | Volume 30 | February 2023 | 216–225
223
Article
https://doi.org/10.1038/s41594-022-00910-8
Encouraged by the examples tested, we defined an automatic
procedure to generate larger models by iterative alignment of pairs
(Methods). We start building all possible dimers in a complex, then
sort them by pDockQ, and start building from the first ranked dimers.
Next, we add the highest-ranked dimer, which shares one subunit with
the complex if it does not overlap; this is repeated for all dimers until
the complex is complete or no additional proteins can be added. We
tested this on the 20S proteasome, a particularly challenging example,
with stoichiometries different from 1-to-1 and homologous subunits.
This automatic procedure could build a model containing all 14 subu-
nits (half of the proteasome), which are mostly placed in agreement
within the experimental model (Fig. 6). However, the exact order of
the chains is incorrect, that is, at each location an incorrect protein is
placed, highlighting that AF2 cannot distinguish which two proteins
interact from a set of homologous proteins.
Two additional proteins where we could build a good model are
Heterodisulfide reductase from Methanothermococcus thermolitho-
trophicus (PDB:5ODC) and the eukaryotic translation initiation factor
2B from Schizosaccharomyces pombe (PDB:5B04) (Supplementary
Fig. 4). For PDB:5ODC we could build a complete model of the protein
with an r.m.s. deviation of 6.0 Å (TM-score 0.90)30 starting from dimers.
However, for PDB:5B04 it was not possible as the chains started over-
lapping when we tried to build a larger model. However, if we build
trimers and then use all three dimers from these trimers we can build a
complete model with an r.m.s. deviation of 7.3 Å (TM-score 0.86), show-
ing that it is sometimes necessary to use larger subunits to assemble
the complexes. Results from a follow-up study31 show that it is often
possible to build the structures of complexes if the subunits are well
predicted. In summary, we find that it is possible to iteratively align
structures of pairs of interacting proteins to build larger assemblies,
but we also identified issues that limit this procedure at the moment.
Concluding discussion
We have predicted complex structures for pairs of human proteins
known to physically interact from two different datasets based on differ-
ent experimental approaches. We note that the source of data used for
the protein interactions is important and impacts the fraction of models
that can be confidently predicted. Our analysis suggests that protein
interactions supported by a combination of affinity-, co-fraction-
and complementation-based methods result in higher-confidence
models. We believe these protein interactions tend to correspond to
high-affinity interactions which are very likely to share a direct physical
permanent interaction. We show that it is possible to use metrics from
the models (for example, pDockQ score) to rank higher-confidence
models, providing an additional accuracy level to large-scale PPI stud-
ies, and in the future to provide additional high-quality targets for
detailed studies of stable complexes. Experimental data from crosslink
mass spectrometry experiments provide an ideal resource for further
validating these predictions via orthogonal means.
Based on comparisons with solved structures, we suggest that
models with pDockQ > 0.5 are 80\% likely to be correct. Additionally,
models with lower scores (pDockQ > 0.23) are still 70\% likely to contain
many correct solutions and may highlight correct interfaces. Such
lower-confidence models are likely to be useful for generating hypoth-
eses and large-scale analyses of global properties. Equally important
is the caveat that high-confidence predictions will still contain errors,
and, in particular, we note that in protein complexes containing paralo-
gous proteins (which is common in higher eukaryotes32), the current
PSMA6
PSMA5
PSMA4
PSMA8
PSMA3
PSMA2
PSMA1
PSMA7
PSMB5
PSMB6
PSMB3
PSMB8
PSMB7
PSMB2
PSMB4
PSMB1
GDI1
RAB22A
RAB12
RAB31
RAB8A
RAB5B
RAB4A
RAB34
RAB4B
RAB13
RAB7A
RAB10
RAB9A
RabGDP
20S proteasome
RFC complex
RFC1
RFC5
RFC3
RFC2
RFC4
RFC3
RFC1
RFC2
RFC4
RFC5
GTF2H4
ERCC3
GTF2H5
ERCC2
GTF2H2C
GTF2H3
GTF2H2
GTF2H1
TFIIH transcription factor core
GTF2H4
GTF2H3
GTF2H1
ERCC3
GTF2H2
Edge/node color
structural evidence
YES
NO
Fig. 6 | Protein complex predictions for higher-order assemblies. The middle
circle is a network view of all PPIs predicted with high confidence (pDockQ > 0.5).
The edges and nodes are colored in red if there is a previous experimental or
homology model for the interaction and blue if such information is unavailable.
We selected four examples of recapitulated complexes (yellow circles and black
arrows) plotted in further detail. In these small networks, only the edges are
colored based on structural evidence. In the case of RabGDP, the faded nodes and
edges represent predictions with slightly lower confidence (pDockQ > 0.3).
Nature Structural \& Molecular Biology | Volume 30 | February 2023 | 216–225
224
Article
https://doi.org/10.1038/s41594-022-00910-8
procedure cannot identify the exact pairing of the protein. For such
cases, additional methods need to be developed.
Structural models for protein interfaces are critical for under-
standing molecular mechanisms and the impact of mutations and
post-translational modifications. We illustrate this using disease muta-
tions and phosphorylation data. While much disease-associated varia-
tion is often found in noncoding regions of the genome, the growth of
exome sequencing of large cohorts of patients will lead to discovering
many more protein mutations linked to disease, which will require such
large structural characteristics. Both for mutations and for phospho-
rylation sites, we think these analyses should be seen as generating
hypotheses for further testing, and we make this information available
in the supplementary material to facilitate such future work.
Finally, we show that it is in principle possible to build structural
models for larger assemblies from predicted binary complexes. In a
follow-up paper we have shown that it is possible to build large assem-
blies fully automatically by using predictions of dimers and trimers31.
Aspects that may limit this include the structural homology between
subunits, unknown subunit stoichiometries and limits in the predicted
interactions31. Additional work will be needed to determine the exact
stoichiometry and to design methods and score systems to build such
larger complex assemblies, as well as to predict the interactions of
proteins with weak and transient interactions.
Online content
Any methods, additional references, Nature Portfolio reporting sum-
maries, source data, extended data, supplementary information,
acknowledgements, peer review information; details of author contri-
butions and competing interests; and statements of data and code avail-
ability are available at https://doi.org/10.1038/s41594-022-00910-8.
References
1.
Orchard, S. et al. The MIntAct project—IntAct as a common
curation platform for 11 molecular interaction databases. Nucleic
Acids Res. 42, D358–D363 (2014).
2.
Luck, K. et al. A reference map of the human binary protein
interactome. Nature 580, 402–408 (2020).
3.
Drew, K., Wallingford, J. B. \& Marcotte, E. M. hu.MAP 2.0:
integration of over 15,000 proteomic experiments builds a global
compendium of human multiprotein assemblies. Mol. Syst. Biol.
17, e10016 (2021).
4.
Mosca, R., Céol, A. \& Aloy, P. Interactome3D: adding structural
details to protein networks. Nat. Methods 10, 47–53 (2012).
5.
Burley, S. K. et al. RCSB Protein Data Bank: powerful new tools
for exploring 3D structures of biological macromolecules for
basic and applied research and education in fundamental
biology, biomedicine, biotechnology, bioengineering and energy
sciences. Nucleic Acids Res. 49, D437–D451 (2021).
6.
Wang, X. et al. Three-dimensional reconstruction of protein
networks provides insight into human genetic disease. Nat.
Biotechnol. 30, 159–164 (2012).
7.
Kamburov, A. et al. Comprehensive assessment of cancer
missense mutation clustering in protein structures. Proc. Natl
Acad. Sci. USA 112, E5486–E5495 (2015).
8.
Porta-Pardo, E., Garcia-Alonso, L., Hrabe, T., Dopazo, J. \& Godzik,
A. A pan-cancer catalogue of cancer driver protein interaction
interfaces. PLoS Comput. Biol. 11, e1004518 (2015).
9.
Beltrao, P. et al. Systematic functional prioritization of protein
posttranslational modifications. Cell 150, 413–425 (2012).
10.	 Nishi, H., Hashimoto, K. \& Panchenko, A. R. Phosphorylation in
protein-protein binding: effect on stability and function. Structure
19, 1807–1815 (2011).
11.
Šoštarić, N. et al. Effects of acetylation and phosphorylation on
subunit interactions in three large eukaryotic complexes. Mol.
Cell. Proteom. 17, 2387–2401 (2018).
12.	 Betts, M. J. et al. Systematic identification of
phosphorylation-mediated protein interaction switches.
PLoS Comput. Biol. 13, e1005462 (2017).
13.	 Zhang, Q. C. et al. Structure-based prediction of protein-
protein interactions on a genome-wide scale. Nature 490,
556–560 (2012).
14.	 Mosca, R., Céol, A., Stein, A., Olivella, R. \& Aloy, P. 3did: a catalog
of domain-based interactions of known three-dimensional
structure. Nucleic Acids Res. 42, D374–D379 (2014).
15.	 Cong, Q., Anishchenko, I., Ovchinnikov, S. \& Baker, D. Protein
interaction networks revealed by proteome coevolution. Science
365, 185–189 (2019).
16.	 Baek, M. et al. Accurate prediction of protein structures and
interactions using a three-track neural network. Science 373,
871–876 (2021).
17.	 Jumper, J. et al. Highly accurate protein structure prediction with
AlphaFold. Nature 596, 583–589 (2021).
18.	 Bryant, P., Pozzati, G. \& Elofsson, A. Improved prediction of
protein-protein interactions using AlphaFold2. Nat. Commun. 13,
1265 (2022).
19.	 Pozzati, G. et al. Limits and potential of combined folding and
docking using PconsDock. Bioinformatics 38, 954–961 (2022).
20.	 Akdel, M. et al. A structural biology community assessment
of AlphaFold 2 applications. Nat. Struct. Mol. Biol. 29,
1056–1067 (2022).
21.	 Evans, R. et al. Protein complex prediction with AlphaFold-
Multimer. Preprint at bioRxiv https://doi.org/10.1101/2021.10.04.
463034 (2021).
22.	 Humphreys, I. R. et al. Computed structures of core eukaryotic
protein complexes. Science 374, eabm4805 (2021).
23.	 Giurgiu, M. et al. CORUM: the comprehensive resource of
mammalian protein complexes—2019. Nucleic Acids Res. 47,
D559–D563 (2019).
24.	 IMEx Consortium Curators et al. Capturing variation impact on
molecular interactions in the IMEx Consortium mutations data
set. Nat. Commun. 10, 10 (2019).
25.	 Delgado, J., Radusky, L. G., Cianferoni, D. \& Serrano, L. FoldX
5.0: working with RNA, small molecules and a new graphical
interface. Bioinformatics 35, 4168–4169 (2019).
26.	 Ochoa, D. et al. The functional landscape of the human
phosphoproteome. Nat. Biotechnol. 38, 365–373 (2020).
27.	 Lawrence, R. T., Searle, B. C., Llovet, A. \& Villén, J. Plug-and-play
analysis of the human phosphoproteome by targeted
high-resolution mass spectrometry. Nat. Methods 13,
431–434 (2016).
28.	 Hornbeck, P. V. et al. PhosphoSitePlus, 2014: mutations, PTMs and
recalibrations. Nucleic Acids Res. 43, D512–D520 (2015).
29.	 Ochoa, D. et al. An atlas of human kinase regulation. Mol. Syst.
Biol. 12, 888 (2016).
30.	 Zhang, Y. \& Skolnick, J. Scoring function for automated
assessment of protein structure template quality. Proteins 57,
702–710 (2004).
31.	 Bryant, P. et al. Predicting the structure of large protein
complexes using AlphaFold and sequential assembly. Nat.
Commun. 13, 6027 (2022).
32.	 Marchant, A. et al. The role of structural pleiotropy and regulatory
evolution in the retention of heteromers of paralogs. eLife 8,
e46754 (2019).
Publisher’s note Springer Nature remains neutral with regard to
jurisdictional claims in published maps and institutional affiliations.
Open Access This article is licensed under a Creative Commons
Attribution 4.0 International License, which permits use, sharing,
adaptation, distribution and reproduction in any medium or format,
Nature Structural \& Molecular Biology | Volume 30 | February 2023 | 216–225
225
Article
https://doi.org/10.1038/s41594-022-00910-8
as long as you give appropriate credit to the original author(s) and the
source, provide a link to the Creative Commons license, and indicate
if changes were made. The images or other third party material in this
article are included in the article’s Creative Commons license, unless
indicated otherwise in a credit line to the material. If material is not
included in the article’s Creative Commons license and your intended
use is not permitted by statutory regulation or exceeds the permitted
use, you will need to obtain permission directly from the copyright
holder. To view a copy of this license, visit http://creativecommons.
org/licenses/by/4.0/.
© The Author(s) 2023
Nature Structural \& Molecular Biology
Article
https://doi.org/10.1038/s41594-022-00910-8
Methods
Protein interaction data and annotations
Human protein pairs known to physically interact were obtained from
the hu.MAP dataset, retaining pairwise protein interactions with ≥0.5
confidence, and most interactions from the HuRI dataset. These inter-
actions were further enriched by obtaining annotations on crosslinked
peptides matched across pairs of interaction proteins, disease-related
mutations and protein phosphorylation sites in the selected proteins.
In addition, all nonhomologous pairs from 12 protein complexes (Sup-
plementary Table 5) and 4,320 protein pairs from 2,102 different pro-
tein complexes (Supplementary Table 6) in CORUM23 were used for
additional analyses. A complete list of all datasets is available from
the supplementary data. A subset of crosslink data was collected from
refs. 33–43, and filtered for peptides unique to only one protein sequence.
A crosslink was considered validated by the structure if the distance
between the epsilon amino groups on the side-chains of the relevant
pair of lysine residues was within 32 Å. Clinical missense variants associ-
ated with disease were collected from ClinVar. We selected only those
having pathogenic or likely pathogenic effects, which were mapped to
Uniprot protein sequences using VarMap. The final list of mutated posi-
tions was then compared with the interface positions. We obtained a list
of protein phosphorylation sites with predicted functional relevance26,
phosphosite annotations28 and regulation of phosphorylation sites
across a large panel of conditions29. These phosphosites were also
mapped to interface positions as defined by the predicted models.
All protein interaction networks were processed using R packages
igraph (v.1.2.5) and qgraph (v.1.9), and further graphical editing was
done using Cytoscape44.
Protein complex prediction
To predict protein complexes of pairwise interactions, we used the
FoldDock pipeline18 based on AlphaFold2 (ref. 17). We used the option
of fused + paired MSAs and ran the model configuration m1-10-1 as this
provides the highest success rate accompanied by a 20-fold speed-up.
Both the fused and paired MSAs were constructed by running HHblits
on every single chain against Uniclust30. The fused MSA was generated
by simply concatenating the output of each of the single-chain HHblits
runs for two interacting chains. The paired MSA was constructed by
combining the top hit for each matching OX identifier between two
interacting chains, using the output from the single-chain HHblits runs.
pDockQ confidence score
To score models, we used features from the predicted complexes to
calculate the predicted DockQ score, pDockQ. This score is defined
with the following sigmoidal equation:
pDockQ =
0.707
1 + e−0.03148(x−388.06) + 0.03138
where
x = average interface plDDT*log(number of interface contacts).
The parameters were optimized to predict the DockQ score using
the dataset from ref. 45. The number of interface contacts is defined as
elsewhere in this paper (any residues with an interface atom within
10 Å of the other chain), and the plDDT is the predicted lDDT score
from AlphaFold2 taken over the interface residues as defined by the
interface contacts.
Building larger complexes from binary protein interactions
A simple procedure to build larger complexes from a set of paired
models was developed. All dimers in the set are by default ranked by
their pDockQ values.
	1.
The building is started from a single dimer, by default the
dimer with the highest pDockQ value. This is referred to as
the ‘complex’.
	2.
All other dimers in the set are then tried to be added to the
‘complex’. Starting with the one with the second highest
pDockQ, a chain is added to the complex if:
	(a)	Exactly one chain of the dimer is identical to one chain in the
complex
	(b)	The structure of these two chains is similar enough (default
TM-score > 0.8)
	(c)	The dimer is then rotated so that the two chains overlap
	(d)	The second chain in the dimer does not clash with more than
25\% of its residues (Cα-Cα distance < 5 Å) with any chain in the
complex.
	3.
If a chain is added, the procedure is started over again and
repeated until no more chains can be added.
Analysis of phosphosites in the protein-protein interfaces
Phosphosite residues in interfaces were identified from a previously
published comprehensive list of known human phosphosites26. Kinases
associated with phosphorylation of interface residues were obtained
from the PhosphositePlus database, and over-representation analysis
of kinases was performed using a hyper-geometric test. Highly regu-
lated interface phosphosites were defined as those with more than
twofold change in phosphorylation in more than two perturbation
conditions across a collated phosphoproteomics dataset comprising
a range of physiological conditions and drug treatments29. Pearson
correlation was calculated amongst these regulated phosphosites and
clusters of co-regulated phosphosites were identified using hierarchi-
cal clustering (‘Ward’ method) of Euclidean distances of the correlation
matrix. Phosphosite clusters were created by cutting the dendrogram
at the appropriate level using the cutree (h = 17) function in R. Phos-
phosite clusters that were significantly regulated in each perturba-
tion condition were identified by a Z-test from the comparison of fold
changes in phosphosite measurements of all phosphosites in a cluster
against the overall distribution of phosphorylation fold changes across
the condition. Gene ontology over-representation of each cluster was
performed separately using a hyper-geometric test in R. The gene
ontology terms were obtained from the c5 category of the Molecular
Signature Database (MSigDBv7.1)46. All over-representation analyses
were performed using the enricher function of the clusterProfiler
package (v.3.12.0)6 in R.
Comparison with other databases
All proteins used here were mapped to UniProt47 to retrieve subcellular
localization, STRING48 for coexpression and other interaction data, and
gtex49 for tissue-specific expression.
Reporting summary
Further information on research design is available in the Nature Port-
folio Reporting Summary linked to this article.
Data availability
All datasets and meta-data are available from https://doi.org/10.17044/
scilifelab.16866202.v1. Further, all models generated as well as some
of the multiple sequence alignments can be found at https://archive.
bioinfo.se/huintaf2/. Source data are provided with this paper.
Code availability
All code used in this project can be found at https://gitlab.com/Elofs-
sonLab/huintaf2/. Tools to run AlphaFold2 for combined folding and
docking can be found at https://gitlab.com/ElofssonLab/FoldDock/.
References
33.	 Yugandhar, K. et al. MaXLinker: proteome-wide cross-link
identifications with high specificity and sensitivity. Mol. Cell.
Proteom. 19, 554–568 (2020).
Nature Structural \& Molecular Biology
Article
https://doi.org/10.1038/s41594-022-00910-8
34.	 Schweppe, D. K. et al. XLinkDB 2.0: integrated, large-scale
structural analysis of protein crosslinking data. Bioinformatics 32,
2716–2718 (2016).
35.	 Klykov, O., van der Zwaan, C., Heck, A. J. R., Meijer, A. B.
\& Scheltema, R. A. Missing regions within the molecular
architecture of human fibrin clots structurally resolved by XL-MS
and integrative structural modeling. Proc. Natl Acad. Sci. USA 117,
1976–1987 (2020).
36.	 Steigenberger, B., Pieters, R. J., Heck, A. J. R. \& Scheltema, R. A.
PhoX: an IMAC-enrichable cross-linking reagent. ACS Cent. Sci. 5,
1514–1522 (2019).
37.	 Klykov, O. et al. Efficient and robust proteome-wide approaches
for cross-linking mass spectrometry. Nat. Protoc. 13, 2964–2990
(2018).
38.	 Fasci, D., van Ingen, H., Scheltema, R. A. \& Heck, A. J. R. Histone
interaction landscapes visualized by crosslinking mass spectrometry
in intact cell nuclei. Mol. Cell. Proteom. 17, 2018–2033 (2018).
39.	 Eliseev, B. et al. Structure of a human cap-dependent 48S
translation pre-initiation complex. Nucleic Acids Res. 46,
2678–2689 (2018).
40.	 Gestaut, D. et al. The chaperonin TRiC/CCT associates with
prefoldin through a conserved electrostatic interface essential for
cellular proteostasis. Cell 177, 751–765.e15 (2019).
41.	 Klatt, F. et al. A precisely positioned MED12 activation helix
stimulates CDK8 kinase activity. Proc. Natl Acad. Sci. USA 117,
2894–2905 (2020).
42.	 Sabath, K. et al. INTS10-INTS13-INTS14 form a functional module
of Integrator that binds nucleic acids and the cleavage module.
Nat. Commun. 11, 3422 (2020).
43.	 Mohamed, W. I. et al. The human GID complex engages two
independent modules for substrate recruitment. EMBO Rep. 22,
e52981 (2021).
44.	 Shannon, P. et al. Cytoscape: a software environment for
integrated models of biomolecular interaction networks. Genome
Res. 13, 2498–2504 (2003).
45.	 Green, A. G. et al. Large-scale discovery of protein interactions at
residue resolution using co-evolution calculated from genomic
sequences. Nat. Commun. 12, 1396 (2021).
46.	 Subramaniam, V., Vincent, I. R. \& Jothy, S. Upregulation and
dephosphorylation of cofilin: modulation by CD44 variant isoform
in human colon cancer cells. Exp. Mol. Pathol. 79, 187–193 (2005).
47.	 The UniProt Consortium. UniProt: the universal protein
knowledgebase. Nucleic Acids Res. 46, 2699 (2018).
48.	 Szklarczyk, D. et al. STRING v11: protein–protein association
networks with increased coverage, supporting functional
discovery in genome-wide experimental datasets. Nucleic Acids
Res. 47, D607–D613 (2019).
49.	 Aguet, F. et al. The GTEx Consortium 
\end{document}
